\begin{abstract}
CLIP-based zero-shot anomaly detection usually scores local regions by their
similarity to normal and abnormal text prototypes. This score is a useful
semantic prior, but it is not sufficient evidence for localization: the same
local response may correspond to a defect in one image and to normal texture,
material granularity, or an object boundary in another. We address this
reference gap through image-conditioned normal reference estimation. Given
frozen CLIP/AnomalyCLIP representations, \method{} selects reliable patch
evidence inside a single test image, estimates an image-conditioned normal
reference, conservatively calibrates the prototypes, and recomputes patch
scores in the same semantic space. The added module performs no test-time
backpropagation, does
not update model parameters, and does not use target-category labels or a
normal reference set. Controlled experiments on MVTec AD and VisA show a
consistent mechanism: direct wavelet-map fusion is a negative control,
GlobalRef remains below the image-conditioned reference, semantic-only
prototype adaptation is a strong baseline, and boundary-aware local reliability
with conservative prototype calibration gives the best controlled row across
all reported metrics.
\end{abstract}
